\section{Discussion and Limitations}\label{sec:discussion}

\paragraph{Why some cases still get through.} The trajectories make it possible to identify the cause of all 40 cases
that bypass \ours{} (13 on \tOne, 27 on \tTwo). Each case falls into exactly one gap of
Proposition~\ref{prop:scope}.

The largest group (25 cases, all confidential-value leaks) occurs because the leaked value sits in a data source
labelled public. The broker only protects values taken from confidential sources, so it does not block these values.
The extractor still recognises their format. The limitation of \rFlow{} therefore lies in how sources are labelled,
not in the extractor. The second group (5 cases) occurs because the destination appears in the user request but is not
an authorised destination. Four cases send to an address that the user merely mentions. One case fills the recipient
field with an ordinary word from the request. The third group (10 cases) alters the content sent to a valid
destination, for example by changing a number. The destination and confidential-value rules do not consider this case,
and only the reminder and the injection detector can block it. On AgentDojo, the bypasses use an operation without a
destination (password change) or a read tool (opening a URL), which are exactly the places the broker does not check.

\tTwo{} was written by the same team after the \tOne{} results were known, so it is not possible to separate the set
being harder from the authors unintentionally targeting weaknesses of the rules. The next evaluation step is an attack
set generated directly from the gaps above.

\paragraph{Language.} Under \cNone, removing diacritics or switching to English does not change the attack success rate
significantly after correction (H4, H5). The confidence intervals still include differences of up to about 12 points,
so the paper does not claim that Vietnamese is more or less vulnerable. The factor clearly affected by language is
phrase-based defence. On Dev, the phrase detector fires on 14 to 19 of 48 standard Vietnamese attacks but on only 3 of
48 English attacks. On AgentDojo, \rSig{} never fires in English but fires 14 times on Vietnamese Slack. \rDest{} and
\rFlow{} are language-independent, whereas \rSig{} is not.

\paragraph{Limitations.} The sealed sets were written by the same team, although the authors of the attacks did not
read the defence code. The data are synthetic, the tools are mocks and the tasks are single-turn. The open models have
only 8 to 12 billion parameters. \ours{} lowers benign utility by about 5 points. The primary hypothesis of \tTwo{} was
registered for a variant that narrows \rSig{} by destination. This variant raises utility by 1.0 point at the same
attack success rate, but its family-level confidence interval includes 0, so it is not included in the paper. Gemini 3
Flash runs only six configurations on Dev and the undefended configuration on AgentDojo. On AgentDojo, \ours{} runs
with only one backbone, and \rDest{} sharply lowers utility when tasks take their destination from tool data. Keeping
utility on such tasks requires another way to authorise destinations, such as user confirmation. The block-all-sinks
reference is an estimate from trajectories. Denial feedback has not been isolated. Extractor coverage over all
trajectories and latency have not been measured.
